\chapter{Machines multi-sockets et interconnexions}\label{ch:ll:multi-socket-machines-and-interconnects}

Un serveur de trading à deux sockets de processeur, une carte réseau dans un emplacement relié au second socket, et un thread de stratégie
que le système d'exploitation a lancé par hasard sur le premier~: chaque paquet reçu par la carte était écrit près du socket 1
et lu depuis le socket 0, et chaque ordre faisait le trajet inverse. Chaque traversée coûte environ une centaine
de nanosecondes de plus que le fait de rester sur un seul socket, plusieurs fois par message, sur un chemin dont le budget logiciel total tient en
quelques microsecondes. Personne n'avait choisi le placement~; personne n'avait regardé la topologie de la machine. Ce chapitre décrit
la machine au-dessus du cœur~: la mémoire rattachée aux sockets, les liens entre sockets, le bus qui relie la carte réseau,
et la manière de lire l'ensemble depuis le système d'exploitation et de placer un processus de trading dessus.

\section{Accès mémoire non uniforme}

\begin{definition}[Accès mémoire non uniforme, nœud NUMA]\label{def:ll:multi-socket-machines-and-interconnects:numa}
Dans une machine à \emph{accès mémoire non uniforme}\index{accès mémoire non uniforme} (NUMA), chaque socket de processeur possède ses
propres contrôleurs mémoire et sa propre mémoire~; un cœur accède plus vite à la mémoire de son propre socket qu'à celle d'un autre socket, laquelle est atteinte
par l'\omterm{def:ll:multi-socket-machines-and-interconnects:link}{interconnexion entre sockets}. Un \emph{nœud NUMA}\index{nœud NUMA} est un ensemble de CPU et de la mémoire qui leur est proche,
en général un socket (ou une partie d'un socket).
\end{definition}

\begin{definition}[Allocation au premier accès]\label{def:ll:multi-socket-machines-and-interconnects:first}
Avec l'\emph{allocation au premier accès}\index{allocation au premier accès} (first touch), la politique par défaut de Linux, une page de mémoire est
placée sur le \omterm{def:ll:multi-socket-machines-and-interconnects:numa}{nœud NUMA} du CPU qui l'écrit en premier (le défaut de page qui la crée), et non sur celui qui a alloué
la plage d'adresses.
\end{definition}

Le premier accès fait du placement une question de savoir qui initialise quoi. Un processus de trading qui construit ses carnets d'ordres, ses files
circulaires et ses réserves dans son thread principal au démarrage, avant de créer ses threads, les place tous sur le nœud du thread principal,
quel que soit le nœud où les threads critiques s'exécutent ensuite. La règle~: épingler d'abord chaque thread critique (chapitre~9), puis le laisser allouer et initialiser
ses propres données~; ou lier la mémoire explicitement (\texttt{numactl --membind}, \texttt{mbind}).

\begin{example}[Un carnet d'ordres alloué sur le mauvais nœud]\label{ex:ll:multi-socket-machines-and-interconnects:touch}
Le thread principal, sur le nœud 0, alloue puis remplit à zéro 64~MiB de carnets d'ordres avec de \omterm{def:ll:memory-hierarchy-and-caches:tlb}{grandes pages} au démarrage, puis lance le
thread du flux de marché sur le nœud 1. Chaque mise à jour du carnet atteint désormais la mémoire de l'autre socket. La correction tient en une ligne~: déplacer
l'initialisation dans le thread du flux de marché, après qu'il s'est épinglé, ou transmettre le nœud à l'allocateur.
\end{example}

\section{L'interconnexion entre sockets et les caches distants}

\begin{definition}[Interconnexion entre sockets]\label{def:ll:multi-socket-machines-and-interconnects:link}
L'\emph{interconnexion entre sockets}\index{interconnexion entre sockets} est l'ensemble des liaisons point à point entre les sockets de processeur
d'une machine, qui acheminent les accès mémoire distants et le trafic de \omterm{def:ll:memory-hierarchy-and-caches:coherence}{cohérence de cache} entre sockets.
\end{definition}

Toute transaction de cohérence du chapitre~3 qui concerne une \omterm{def:ll:memory-hierarchy-and-caches:line}{ligne de cache} détenue sur l'autre socket traverse l'interconnexion. Une
mesure publiée sur un serveur à deux sockets d'une génération récente d'Intel a montré que le rebond d'une \omterm{def:ll:memory-hierarchy-and-caches:line}{ligne de cache} entre deux
cœurs prenait en moyenne 59~nanosecondes au sein d'un même socket (81 au pire) et 138~nanosecondes d'un socket à l'autre~; les
anciennes générations étaient similaires, jusqu'à 150 d'un socket à l'autre. Un passage de relais entre deux threads, le feed handler publiant un message et
la stratégie qui le lit, est exactement un tel rebond (chapitre~12)~: placé sur deux sockets, son coût plus que double.

La même mesure sur l'ordinateur portable de ce livre, qui possède un seul socket et tourne sous un hyperviseur, est
la \cref{fig:ll:multi-socket-machines-and-interconnects:c2c}. Elle présente une structure~: des paires de CPU virtuels voisins se transmettent une
\omterm{def:ll:memory-hierarchy-and-caches:line}{ligne de cache} en environ 40~nanosecondes, la plupart des autres paires en 170 à 200. Les cœurs physiques de l'ordinateur portable sont regroupés en grappes partageant
un cache de second niveau, et les CPU virtuels leur sont associés par l'hyperviseur, sans que le système invité le voie~; la
matrice, et non la numérotation des CPU, qui indique quels threads doivent communiquer entre eux.

\begin{omfigure}
\begin{tikzpicture}
\begin{axis}[width=8.6cm,height=8.6cm,view={0}{90},xmin=-0.5,xmax=21.5,ymin=-0.5,ymax=21.5,y dir=reverse,
  xtick={0,3,6,9,12,15,18,21},ytick={0,3,6,9,12,15,18,21},xlabel={CPU},ylabel={CPU},
  tick label style={font=\small},label style={font=\small},
  colormap={om}{color(0)=(white) color(0.3)=(omDef!30) color(1)=(omDef)},
  colorbar,colorbar style={ylabel={ns, sens unique},ylabel style={font=\small},tick label style={font=\footnotesize}},
  point meta min=0,point meta max=200,table/col sep=comma]
\addplot[matrix plot*,mesh/cols=22,point meta=explicit] table[x=b,y=a,meta=ns]{figdata/low-latency/04-multi-socket-machines-and-interconnects/measured_c2c.csv};
\end{axis}
\end{tikzpicture}

\omcaption{Latence, dans un sens, du passage d'une \omterm{def:ll:memory-hierarchy-and-caches:line}{ligne de cache} entre deux threads épinglés sur deux des 22 CPU virtuels de l'ordinateur portable
(20\,000 allers-retours, meilleur de trois). Les CPU voisins sont rapides~; les autres sont quatre à cinq fois plus lents. Mesuré sur un
ordinateur portable (Intel Core Ultra 7 155H) sous WSL2, sans cœurs isolés. Données~: \texttt{bench\_c2c.py}.}\label{fig:ll:multi-socket-machines-and-interconnects:c2c}
\end{omfigure}

\section{Le bus périphérique et l'accès direct à la mémoire}

\begin{definition}[PCI Express, accès direct à la mémoire]\label{def:ll:multi-socket-machines-and-interconnects:pcie}
\emph{PCI Express}\index{PCI Express} (PCIe) est le bus série point à point qui relie les périphériques (cartes réseau,
stockage, accélérateurs) à un socket de processeur, en lignes regroupées par 1 à 16.
Dans l'\emph{accès direct à la mémoire}\index{accès direct à la mémoire} (DMA), un périphérique lit et écrit lui-même la mémoire de l'hôte, sans que le processeur recopie chaque octet~: une carte réseau
écrit chaque paquet reçu dans un tampon que le pilote lui a fourni et signale son arrivée par un descripteur, une interruption
ou par un indicateur que le programme interroge.
\end{definition}

\begin{definition}[Accès direct au cache]\label{def:ll:multi-socket-machines-and-interconnects:dca}
L'\emph{accès direct au cache}\index{accès direct au cache} (DCA) permet aux écritures DMA d'un périphérique d'atterrir dans le cache de dernier niveau
plutôt que dans la mémoire, afin que le cœur qui traite les données les trouve dans le cache~; l'implémentation d'Intel (Data Direct I/O)
est transparente pour le logiciel.
\end{definition}

Chaque emplacement PCIe appartient à un socket~: ses lignes proviennent du processeur de ce socket. Les paquets d'une carte arrivent donc dans
le cache ou la mémoire de ce socket, et les registres de la carte, que le programme écrit pour émettre, sont de l'autre côté de l'interconnexion
depuis n'importe quel cœur de l'autre socket. Chaque transaction sur le bus ajoute sa propre latence, raison pour laquelle les cartes réseau à faible latence et leurs pilotes en espace utilisateur
réduisent le nombre de transactions sur le bus par paquet (le livre~14, chapitre~3).

\begin{dated}{2026-09}{Générations de PCI Express}\label{dat:ll:multi-socket-machines-and-interconnects:pcie}
PCIe 4.0 transmet à 16~GT/s par ligne et PCIe 5.0, dont la spécification publiée par le PCI-SIG en mai 2019, à 32~GT/s~; les deux
encodent 128 bits dans 130. Une liaison PCIe 5.0 à 16 lignes transporte donc $32 \times 16 \times 128/130/8 \approx 63$~Go/s dans chaque
sens, bien plus que n'en exige un port réseau de 100~Gb/s (12.5~Go/s)~: pour une carte de trading, le bus compte pour sa latence,
et non sa bande passante.
\end{dated}

\section{Où se place la carte réseau}

\begin{definition}[Localité du périphérique]\label{def:ll:multi-socket-machines-and-interconnects:locality}
La \emph{localité du périphérique}\index{localité du périphérique} d'un thread ou d'un tampon est son \omterm{def:ll:multi-socket-machines-and-interconnects:numa}{nœud NUMA} par rapport au périphérique avec lequel il
échange des données~: locale lorsqu'ils se trouvent sur le même nœud, distante sinon.
\end{definition}

\begin{omfigure}
\begin{tikzpicture}[font=\footnotesize,
  sock/.style={draw=omDef,fill=omDef!6,rounded corners=3pt,minimum width=4.2cm,minimum height=2.4cm},
  mem/.style={draw=gray,fill=gray!12,minimum width=4.2cm,minimum height=0.55cm},
  thr/.style={draw=omThm,fill=omThm!12,rounded corners=2pt,minimum width=1.6cm,minimum height=0.5cm}]
\node[mem] (m0) at (0,2.1) {mémoire du nœud 0};
\node[mem] (m1) at (7,2.1) {mémoire du nœud 1};
\node[sock] (s0) at (0,0) {};
\node[sock] (s1) at (7,0) {};
\node[anchor=north] at (s0.north) {socket 0~: cœurs, caches};
\node[anchor=north] at (s1.north) {socket 1~: cœurs, caches};
\node[thr] (bad) at (0,-0.35) {stratégie (distante)};
\node[thr,fill=omMeth!15,draw=omMeth] (good) at (7,-0.35) {stratégie (locale)};
\draw[thick] (m0) -- (s0); \draw[thick] (m1) -- (s1);
\draw[<->,very thick,omProp] (s0.east) -- node[above,align=center]{interconnexion\\ 138 ns par ligne} (s1.west);
\node[draw=black,fill=gray!20,minimum width=2.4cm,minimum height=0.6cm] (nic) at (7,-2.3) {carte réseau};
\draw[<->,thick] (nic) -- node[right]{lignes PCIe du socket 1} (s1.south);
\end{tikzpicture}

\omcaption{Un serveur à deux sockets dont la carte réseau est reliée au socket 1. Un thread de stratégie sur le socket 1 lit les paquets depuis
son propre cache~; un thread sur le socket 0 va chercher chaque ligne de paquet, et envoie chaque ordre, à travers l'interconnexion, à environ 138
nanosecondes par ligne, contre 59 à l'intérieur d'un socket (moyennes publiées pour une génération récente).}\label{fig:ll:multi-socket-machines-and-interconnects:server}
\end{omfigure}

\begin{proposition}[Coût d'un placement distant]\label{prop:ll:multi-socket-machines-and-interconnects:remote}
Si un message apporte $k_{\mathrm{in}}$ \omterm{def:ll:memory-hierarchy-and-caches:line}{lignes de cache} de la carte vers le thread critique, et qu'un ordre en renvoie $k_{\mathrm{out}}$,
chacune transférée l'une après l'autre, et qu'une ligne coûte $c_{\mathrm L}$ à l'intérieur du socket et $c_{\mathrm R}$ à travers lui, alors
le placement distant ajoute $(k_{\mathrm{in}} + k_{\mathrm{out}})(c_{\mathrm R} - c_{\mathrm L})$ à la \omterm{def:ll:where-latency-comes-from:t2t}{latence tick-to-trade}.
\end{proposition}

\begin{proof}
Chacun des $k_{\mathrm{in}} + k_{\mathrm{out}}$ transferts dépendants coûte $c_{\mathrm R}$ au lieu de $c_{\mathrm L}$~; le reste
du chemin reste inchangé.
\end{proof}

Avec deux lignes dans chaque sens et les valeurs publiées de 59 et 138 nanosecondes, la proposition donne 316 nanosecondes par message~:
plus d'un dixième des \qty{2.5}{\micro\second} qu'un très bon système logiciel met du fil à la sortie du fil (chapitre~1). Les
règles de placement en découlent~: la carte, les threads critiques et leur mémoire sur un même nœud~; chaque thread critique sur son propre cœur physique,
son frère d'hyper-threading inactif~; les interruptions des autres périphériques et le travail du système d'exploitation sur l'autre nœud ou
sur des cœurs de service (chapitre~13)~; et chaque thread qui échange des données avec le chemin critique (hot path) par message (la barrière de risque (risk gate),
l'anneau du journal) sur le même nœud.

\section{Lire la topologie d'une machine}

Linux expose la topologie dans sysfs. Le boîtier physique, le cœur et les frères d'hyper-threading de chaque CPU se trouvent sous
\path{/sys/devices/system/cpu/cpu*/topology/}~; les CPU de chaque nœud se trouvent dans \path{/sys/devices/system/node/node*/cpulist}~;
le nœud de chaque périphérique PCI se trouve dans \path{/sys/bus/pci/devices/*/numa_node}, qui contient $-1$ lorsque le firmware ne l'a pas indiqué~; et
le périphérique derrière une interface réseau est le lien \path{/sys/class/net/<name>/device}. Des outils tels que
\texttt{lscpu}, \texttt{numactl --hardware} et \texttt{lstopo} affichent les mêmes informations~; un plan de placement doit les lire
lui-même, afin qu'une carte déplacée ou un nouveau serveur ne puisse pas invalider le plan en silence.

Cet ordinateur portable montre ce qui peut mal tourner. Linux voit un seul nœud, 22 CPU décrits comme onze cœurs à deux threads (le processeur
compte 6 cœurs de performance à deux threads et 10 cœurs efficaces à un seul thread), et une interface réseau qui est un adaptateur
synthétique de l'hyperviseur, et non un périphérique PCI~: il n'existe aucune carte dont il faudrait être local. La construction de ce chapitre
traite donc un nœud inconnu comme un fait à signaler, et non comme une valeur à deviner.

\section{Tutoriel~: un plan de placement et la matrice cœur à cœur}\label{tut:ll:multi-socket-machines-and-interconnects:tut}

\begin{tutorial}
\textbf{Objectif.} Lire la topologie d'une machine, placer les threads critiques d'un processus de trading sur celle-ci, vérifier le plan, et mesurer
le coût d'un passage de relais entre chaque paire de CPU. \textbf{État final~:} un plan vérifié pour une maquette à deux sockets et pour
cet ordinateur portable, et la \cref{fig:ll:multi-socket-machines-and-interconnects:c2c}.
\begin{tutsteps}
\item \textbf{Lire l'arborescence.} \texttt{read\_topology(root)} analyse les fichiers sysfs décrits plus haut~; les tests l'exécutent sur deux
arborescences synthétiques écrites par \texttt{make\_coreplan\_fixture.py} (un serveur à deux sockets de 32 CPU dont la carte est sur le nœud 1, et une
machine à un seul nœud dont le nœud de la carte est inconnu).
\item \textbf{Placer.} Les rôles critiques prennent les cœurs physiques de numéro le plus élevé du nœud de la carte, un chacun, en laissant leurs
frères inactifs~; le cœur du CPU 0 et tous les autres CPU sont confiés au service.
\omcode{code/firm/coreplan/firm_coreplan.py}{115}{131}{Un plan de placement à partir de la topologie.}\label{lst:ll:multi-socket-machines-and-interconnects:plan}
\item \textbf{Vérifier.} \texttt{check(plan, topo)} liste chaque règle violée~: un thread critique sur le mauvais nœud, deux threads critiques
sur un même cœur physique, un frère actif, le CPU 0 hors du service.
\item \textbf{Mesurer les passages de relais.} Deux threads épinglés se renvoient un drapeau atomique rembourré~; le temps par transfert
dans un sens est la latence cœur à cœur. \texttt{python bench\_c2c.py} l'exécute pour chaque paire.
\omcode{code/low-latency/04-multi-socket-machines-and-interconnects/cpp/ll_c2c.hpp}{23}{50}{Renvoi d'une ligne de cache entre deux threads épinglés.}\label{lst:ll:multi-socket-machines-and-interconnects:pingpong}
\end{tutsteps}
\textbf{Pour aller plus loin.} Ajouter une règle à \texttt{check} pour que le journal soit lui aussi sur le nœud de la carte, et décider
s'il doit l'être~; relancer le ping-pong avec le drapeau et une charge de quatre lignes, puis comparer le coût par ligne.
\end{tutorial}

\section{Construction~: le plan de placement des cœurs}\label{bld:ll:multi-socket-machines-and-interconnects:coreplan}

\begin{build}
\bfield{Rôle} Une description unique de l'endroit où s'exécute chaque thread du processus de trading, dérivée de la machine plutôt que
saisie à la main~: le chapitre~9 épingle les threads à partir d'elle, le chapitre~13 audite un hôte par rapport à elle, le chapitre~26 l'enregistre avec chaque
mesure.
\bfield{Interface} Python \texttt{firm\_coreplan}~: \texttt{read\_topology(root)} renvoyant \texttt{Topology} (CPU avec
boîtier, cœur, nœud et frères~; nœuds~; cartes PCI avec nœud et CPU locaux~; interfaces), \texttt{Topology.card\_for(ifname)},
\texttt{plan(topo, ifname, hot, cold)} renvoyant \texttt{Plan(node, roles, idle, housekeeping)},
\texttt{check(plan, topo)}, \texttt{remote\_penalty\_ns}.
\bfield{Règles} Rôles critiques sur le nœud de la carte, un cœur physique chacun, frères inactifs, jamais le cœur du CPU 0~; un nœud inconnu
($-1$) se rabat sur le nœud le plus bas et est signalé~; le plan échoue bruyamment lorsque le nœud a trop peu de cœurs.
\bfield{Tests d'acceptation} \texttt{code/firm/coreplan/tests/}~: analyse des listes de CPU~; la topologie de la maquette à deux sockets et un
plan sûr sur le nœud 1~; les plans dont un thread est sur le mauvais socket et ceux dont deux threads sont sur un même cœur sont détectés~; la maquette à un seul nœud avec
une carte inconnue.
\bfield{Pour aller plus loin} Lire la topologie du cache (\path{/sys/devices/system/cpu/cpu*/cache/}) et placer les threads qui échangent des données par
message sur des cœurs partageant un niveau de cache~; inclure la matrice cœur à cœur mesurée dans le plan.
\end{build}

\begin{omsources}
\item Chips and Cheese, ``Core to core latency data on large systems'', 2023.
\item Documentation du noyau Linux, ``NUMA memory policy''~; documentation ABI, \texttt{sysfs-bus-pci}.
\item Intel, ``Intel Data Direct I/O technology''.
\item PCI-SIG, PCI Express 5.0 specification announcement, 2019.
\end{omsources}

\section{Exercices}

\begin{exercise}[$\star$]\label{exo:ll:multi-socket-machines-and-interconnects:1}
Quelle est la bande passante utile dans chaque sens d'un lien PCIe 4.0 à 8 voies~? Suffit-elle pour deux ports de 25~Gb/s~?
\end{exercise}

\begin{exercise}[$\star$]\label{exo:ll:multi-socket-machines-and-interconnects:2}
Un paquet de 300 octets arrive dans le cache du socket de la carte. Combien de \omterm{def:ll:memory-hierarchy-and-caches:line}{lignes de cache} un cœur de l'autre socket doit-il aller chercher
pour le lire~?
\end{exercise}

\begin{exercise}[$\star$]\label{exo:ll:multi-socket-machines-and-interconnects:3}
\path{/sys/bus/pci/devices/0000:81:00.0/numa_node} contient 1 et la stratégie s'exécute sur le CPU 5, dont le nœud est 0. Que
donne la \cref{prop:ll:multi-socket-machines-and-interconnects:remote} avec trois lignes entrantes, une sortante, et les
valeurs publiées~?
\end{exercise}

\begin{exercise}[$\star\star$]\label{exo:ll:multi-socket-machines-and-interconnects:4}
Le feed handler, épinglé sur le nœud 1, publie chaque message dans un anneau que le thread principal a alloué et mis à zéro au
démarrage sur le nœud 0. Où va chaque écriture, et comment corriger cela~?
\end{exercise}

\begin{exercise}[$\star\star$]\label{exo:ll:multi-socket-machines-and-interconnects:5}
Lire la \cref{fig:ll:multi-socket-machines-and-interconnects:c2c}~: quelles paires de CPU attribuer à un thread de flux et à un
thread de stratégie qui se passent chaque message, et quel serait le coût d'un passage de relais entre la paire choisie et une paire quelconque~?
\end{exercise}

\begin{exercise}[$\star\star$]\label{exo:ll:multi-socket-machines-and-interconnects:6}
Expliquer pourquoi le frère hyper-threading d'un thread chaud doit rester inactif, et ce que le plan y perd sur l'installation à deux sockets~?
\end{exercise}

\begin{exercise}[$\star\star\star$]\label{exo:ll:multi-socket-machines-and-interconnects:7}
\emph{Programmation.} Étendre \texttt{firm\_coreplan.plan} pour placer quatre threads chauds sur un nœud ne disposant que de trois cœurs physiques utilisables
en partageant un cœur entre les deux threads qui échangent le moins de données. Quelle paire choisir dans un pipeline de flux, de carnet,
de stratégie et de passerelle, et que doit accepter \texttt{check}~?
\end{exercise}

\begin{exercise}[$\star\star\star$]\label{exo:ll:multi-socket-machines-and-interconnects:8}
\emph{Trouver la faille.} «~Nous avons épinglé notre stratégie avec \texttt{taskset -c 5} dans le script de démarrage du service, elle est
donc sur le bon socket.~»
\end{exercise}

\section{Problème~: mauvais socket}

\begin{problem}[{Problème du week-end --- un serveur que personne n'avait examiné}]\label{pb:ll:multi-socket-machines-and-interconnects:1}
Le nouveau serveur à deux sockets d'une société a sa carte réseau sur le nœud 1. Le feed handler, la stratégie et la passerelle ont été lancés par un
script sans aucun épinglage, et l'ordonnanceur les a placés sur le nœud 0. Chaque message de données de marché apporte deux \omterm{def:ll:memory-hierarchy-and-caches:line}{lignes de cache} à la
stratégie~; chaque ordre envoie deux lignes à la carte. Prendre les chiffres publiés~: 59 nanosecondes par ligne au sein d'un socket, 138
entre sockets.

\medskip\noindent\textbf{Partie I --- La topologie.}
\begin{enumerate}
\item Quels fichiers sysfs indiquent le nœud de la carte et celui de chaque CPU~?
\item Sur l'installation \texttt{two\_socket}, quels CPU se trouvent sur le nœud 1, et quels sont les frères du CPU 17~?
\item Que renvoie \texttt{plan} pour les trois rôles chauds, et quels CPU laisse-t-il inactifs~?
\item Pourquoi le plan n'utilise-t-il jamais le cœur du CPU 0~?
\end{enumerate}

\medskip\noindent\textbf{Partie II --- La pénalité.}
\begin{enumerate}[resume]
\item Appliquer la \cref{prop:ll:multi-socket-machines-and-interconnects:remote}~: que coûte par message le mauvais socket~?
\item Quelle fraction d'un budget fil-à-fil de \qty{2.5}{\micro\second} cela représente-t-elle~?
\item Si la stratégie lit aussi son carnet d'ordres, alloué par le thread principal sur le nœud 0, que se passe-t-il lorsque les threads sont
déplacés sur le nœud 1 sans déplacer le carnet~?
\item Laquelle des deux erreurs est la plus grave, et pourquoi~?
\end{enumerate}

\medskip\noindent\textbf{Partie III --- Passages de relais.}
\begin{enumerate}[resume]
\item Le thread de flux remet chaque message à la stratégie par un anneau~: une ligne dans chaque sens. Quel est le coût du passage de relais
au sein d'un socket, et entre sockets~?
\item Sur cet ordinateur portable, combien coûte un passage entre CPU voisins, et entre les autres (\cref{fig:ll:multi-socket-machines-and-interconnects:c2c})~?
\item Comment choisir les deux CPU à partir de la seule matrice~?
\item Pourquoi la matrice peut-elle changer d'un démarrage de machine virtuelle à l'autre~?
\end{enumerate}

\medskip\noindent\textbf{Partie IV --- La correction.}
\begin{enumerate}[resume]
\item Énoncer le \emph{résultat nommé}~: la latence ajoutée par message par le mauvais socket, en nanosecondes et en fraction
d'un budget de \qty{2.5}{\micro\second}.
\item Écrire les étapes de la correction, dans l'ordre.
\item Comment vérifier, après la correction, que chaque thread s'exécute là où le plan le prévoit~?
\item Que doit-il se passer lorsqu'un changement matériel déplace la carte sur le nœud 0~?
\item Les interruptions d'autres périphériques doivent-elles être éloignées des cœurs chauds~?
\item Que signifie \texttt{numa\_node} $= -1$, et que doit faire le plan~?
\item Que faut-il consigner avec chaque mesure de latence pour qu'un changement de placement reste visible plus tard~?
\item En une phrase~: pourquoi le placement est-il une décision de latence~?
\end{enumerate}
\end{problem}

\section{Questions d'entretien}

\begin{interviewq}[$\star$ \iqroles{développeur}]\label{iq:ll:multi-socket-machines-and-interconnects:1}
Qu'est-ce que NUMA, et que signifie-t-il pour un processus sensible à la latence~?
\end{interviewq}

\begin{interviewq}[$\star\star$ \iqroles{développeur}]\label{iq:ll:multi-socket-machines-and-interconnects:2}
Qu'est-ce que l'\omterm{def:ll:multi-socket-machines-and-interconnects:first}{allocation au premier accès} (first touch), et comment peut-elle placer la mémoire d'un thread sur le mauvais nœud~?
\end{interviewq}

\begin{interviewq}[$\star\star$ \iqroles{développeur}]\label{iq:ll:multi-socket-machines-and-interconnects:3}
Comment savoir, sous Linux, à quel socket est reliée une carte réseau~?
\end{interviewq}

\begin{interviewq}[$\star\star$ \iqroles{développeur}]\label{iq:ll:multi-socket-machines-and-interconnects:4}
Deux threads échangent un message par microseconde. Où les placeriez-vous sur une machine à deux sockets, et pourquoi~?
\end{interviewq}

\begin{interviewq}[$\star\star$ \iqroles{développeur}]\label{iq:ll:multi-socket-machines-and-interconnects:5}
Qu'est-ce que le DMA, et qu'est-ce que l'\omterm{def:ll:multi-socket-machines-and-interconnects:dca}{accès direct au cache} change pour le chemin d'un paquet jusqu'à l'application~?
\end{interviewq}

\begin{interviewq}[$\star\star\star$ \iqroles{développeur}]\label{iq:ll:multi-socket-machines-and-interconnects:6}
Concevez une procédure de démarrage pour un processus de trading qui garantisse le bon placement de ses threads et de sa mémoire sur
n'importe quel serveur de la flotte.
\end{interviewq}
